\section*{Chapter \ref{ch:rs:signal-combination} --- Signal Combination}

\begin{solution}{exo:rs:signal-combination:1}
$10 \times 0.02/\sqrt{10 + 90 \times 0.25} = 0.2/\sqrt{32.5} = 0.035$; the limit is $0.02/\sqrt{0.25} = 0.04$.
\end{solution}

\begin{solution}{exo:rs:signal-combination:2}
IC weights use the positive parts 0.03, 0.01 and 0: weights 0.75, 0.25 and 0. Halfway toward equal weights:
$\frac12(0.75, 0.25, 0) + \frac12(\frac13, \frac13, \frac13) = (0.542, 0.292, 0.167)$.
\end{solution}

\begin{solution}{exo:rs:signal-combination:3}
Returns and signals are ordered in time, correlated across neighbouring dates and exposed to regimes; a fold that trains on later dates and tests on
earlier ones lets the \omterm{def:rs:signal-combination:stacking}{stacking} weights learn from the future (which blend worked in the test period's own regime). Expanding windows that always
predict forward reproduce what the firm can do in production.
\end{solution}

\begin{solution}{exo:rs:signal-combination:4}
$w \propto C^{-1}\rho$ with $C = \begin{pmatrix} 1 & 0.5\\ 0.5 & 1\end{pmatrix}$: $C^{-1}\rho = \frac{1}{0.75}(0.03 - 0.01, -0.015 + 0.02) = (0.0267, 0.0067)$,
weights 0.8 and 0.2. The best IC is $\sqrt{\rho^\top C^{-1}\rho} = \sqrt{0.03 \times 0.0267 + 0.02 \times 0.0067} = 0.0306$; equal weights give
$0.05/\sqrt{2 + 2 \times 0.5} = 0.0289$.
\end{solution}

\begin{solution}{exo:rs:signal-combination:5}
The first signal is its own first Gram--Schmidt vector. The second, residualised on the first, is $z_2 - cz_1$, with variance $1 - c^2$; standardised,
$(z_2 - cz_1)/\sqrt{1 - c^2}$. Its correlation with the target is $(\rho_2 - c\rho_1)/\sqrt{1 - c^2}$: what the second signal knows that the first
does not.
\end{solution}

\begin{solution}{exo:rs:signal-combination:6}
Sixty observations of forty series give a sample covariance whose smallest eigenvalues are far below the true ones (the spread of sample eigenvalues
of Book~4, chapter~22), and $\Sigma^{-1}\mu$ divides by them: directions that look almost riskless because of sampling error receive the largest
weights. Shrinking toward the diagonal lifts those eigenvalues; at the diagonal the rule becomes mean over variance per signal, close to IC weights.
\end{solution}

\begin{solution}{exo:rs:signal-combination:7}
The optimal $\lambda$ (on a grid of tenths) is 0.5 at 6 months (IC 0.074), 0.4 at 12, 0.6 at 24 and 0.4 at 48 (0.083). It does not move steadily toward
the regression weights: with 100 names even four years leave enough error in forty regression weights that a large pull toward equal weights pays, and
the optimum is noisy from one window to the next.
\end{solution}

\begin{solution}{exo:rs:signal-combination:8}
The weights were fitted on the period they are judged on: the comparison measures the fit, not the forecast. Refit on dates before each judged date
(expanding windows or time-ordered folds), compare out of sample, and report equal weights alongside; on the chapter's full panel the in-sample
advantage of the unshrunk maximum-ICIR rule reversed out of sample in the drifting world.
\end{solution}

\begin{solution}{pb:rs:signal-combination:1}
\begin{enumerate}
\item 0.015 and 0.027.
\item 0.065 from the proposition; 0.064 measured.
\item $\bar\rho/\sqrt{\bar c} = 0.090$. Each new signal repeats part of its theme's noise, shared with the other signals of the theme; that shared part does
not average away, and the blend's IC approaches the ceiling only as the number of signals grows without bound.
\item Five signals per theme share the theme's noise.
\item Stable world: 0.064, 0.099 and 0.102. Drifting world: 0.060, 0.075 and 0.087.
\item It inverts the covariance of forty IC series estimated from 59 months and loads on its smallest, most underestimated eigenvalues.
\item 0.098 in the stable world and 0.074 in the drifting one: close to the IC-weighted result.
\item The themes' information drifts with a half-life of about two years, so the first five years still say something about the next ones, and 704 names
estimate the weights precisely.
\item 0.073 in the stable world, 0.005 in the drifting one.
\item \textbf{Named result.} With 100 names and a year of history, the best out-of-sample blend in the stable world is IC weights (0.086 against 0.069
for equal weights), and regression weights are best shrunk 40\% toward equal weights (0.081); in the drifting world the best blend is equal weights
($\lambda = 1$, 0.069), and unshrunk regression weights earn 0.005.
\item Stable world: 0.085, 0.086, 0.082 and 0.092 against 0.069, better at every window. Drifting world: 0.043, 0.049, 0.067 and 0.074, worse until four
years of history.
\item The finite-sample error in estimating the combining weights.
\item All its weight on IC weights, in both worlds. It chooses among its base blends on held-out dates; it cannot be better than the best of them by more
than their errors' diversification, and here one base dominated.
\item 0.064 for the symmetric set, as for the raw signals; 0.050 for the sequential set. The sequential order gives the best signal its theme's shared
part and leaves the others with residual noise, which equal weights then overweight; the symmetric set changes each signal as little as possible.
\item When weights must be interpreted (marginal information), when the book's exposures must be removed (residualise on the risk model), or when a new
signal must be judged on what it adds; not to create information.
\item On time-ordered folds within the training period, choosing the $\lambda$ that maximises the average out-of-fold IC, refitted on a schedule, with
equal weights as the default when the choice is not clearly better.
\item Signals of very different and persistent quality, or signals with negative ICs: equal weights keep them all.
\item The panel is large (704 names a month), so the unpenalised regression is already precise, and penalising toward zero mostly rescales a composite
that is standardised anyway.
\item The \omterm{def:rs:signal-combination:blend}{equal-weight blend}'s IC on the same dates, in sample and out of sample.
\item When signal quality differs persistently and the sample is large enough that the weights' estimation error is smaller than the differences they
estimate.
\end{enumerate}
\end{solution}

\begin{solution}{iq:rs:signal-combination:1}
Standardise each by date, start from equal weights, and measure the blend's IC out of sample; then try IC weights and regression weights shrunk toward
equal weights, with the shrinkage chosen on time-ordered folds; keep the estimated weights only if they beat equal weights out of sample by more than
their noise. Check the signals' correlation structure first: the blend's ceiling is set by it.
\end{solution}

\begin{solution}{iq:rs:signal-combination:2}
Simple averages of forecasts repeatedly beat combinations with estimated weights (Stock and Watson; Smith and Wallis). The explanation: the error in
estimating the weights from finite samples, and the instability of the forecasts' relative quality, can exceed the gain from weighting.
\end{solution}

\begin{solution}{iq:rs:signal-combination:3}
Build base blends on expanding windows, predict the following block of dates with each, and fit non-negative \omterm{def:rs:signal-combination:stacking}{stacking} weights on those out-of-sample
predictions only; never let a base or the stacker see dates it is later judged on, and compare the stack with its best base and with equal weights.
\end{solution}

\begin{solution}{iq:rs:signal-combination:4}
It may repeat information already in the composite (high correlation with existing signals: its residual IC, $(\rho - c\rho_1)/\sqrt{1 - c^2}$ for one
correlated signal, can be near zero); it may be noisier than its IC suggests; or the weighting may have given it too little weight to matter. Measure its
incremental IC after residualising on the composite (chapter~12).
\end{solution}

\begin{solution}{iq:rs:signal-combination:5}
Sequential (Gram--Schmidt) \omterm{def:rs:signal-combination:orth}{orthogonalisation} keeps the first signal and residualises each next one on those before: the order decides who keeps the shared
information, useful when there is a natural priority (a book's factors first, new signals after). \omterm{def:rs:signal-combination:orth}{Symmetric orthogonalisation} ($ZS^{-1/2}$) treats all
signals alike and gives the uncorrelated set closest to the originals: useful for interpretation without an order.
\end{solution}

\begin{solution}{iq:rs:signal-combination:6}
Maximise $w^\top\rho/\sqrt{w^\top Cw}$: the gradient condition gives $\rho(w^\top Cw) = Cw(w^\top\rho)$, so $w \propto C^{-1}\rho$ and the maximum is
$\sqrt{\rho^\top C^{-1}\rho}$. Estimated $C^{-1}$ amplifies the sampling error of $C$'s smallest eigenvalues and $\hat\rho$ is itself noisy, so the
weights are extreme and unstable; shrink $C$ toward its diagonal and $w$ toward equal weights.
\end{solution}
